% generated by GAPDoc2LaTeX from XML source (Frank Luebeck)
\documentclass[a4paper,11pt]{report}

\usepackage[top=37mm,bottom=37mm,left=27mm,right=27mm]{geometry}
\sloppy
\pagestyle{myheadings}
\usepackage{amssymb}
\usepackage[utf8]{inputenc}
\usepackage{makeidx}
\makeindex
\usepackage{color}
\definecolor{FireBrick}{rgb}{0.5812,0.0074,0.0083}
\definecolor{RoyalBlue}{rgb}{0.0236,0.0894,0.6179}
\definecolor{RoyalGreen}{rgb}{0.0236,0.6179,0.0894}
\definecolor{RoyalRed}{rgb}{0.6179,0.0236,0.0894}
\definecolor{LightBlue}{rgb}{0.8544,0.9511,1.0000}
\definecolor{Black}{rgb}{0.0,0.0,0.0}

\definecolor{linkColor}{rgb}{0.0,0.0,0.554}
\definecolor{citeColor}{rgb}{0.0,0.0,0.554}
\definecolor{fileColor}{rgb}{0.0,0.0,0.554}
\definecolor{urlColor}{rgb}{0.0,0.0,0.554}
\definecolor{promptColor}{rgb}{0.0,0.0,0.589}
\definecolor{brkpromptColor}{rgb}{0.589,0.0,0.0}
\definecolor{gapinputColor}{rgb}{0.589,0.0,0.0}
\definecolor{gapoutputColor}{rgb}{0.0,0.0,0.0}

%%  for a long time these were red and blue by default,
%%  now black, but keep variables to overwrite
\definecolor{FuncColor}{rgb}{0.0,0.0,0.0}
%% strange name because of pdflatex bug:
\definecolor{Chapter }{rgb}{0.0,0.0,0.0}
\definecolor{DarkOlive}{rgb}{0.1047,0.2412,0.0064}


\usepackage{fancyvrb}

\usepackage{mathptmx,helvet}
\usepackage[T1]{fontenc}
\usepackage{textcomp}


\usepackage[
            pdftex=true,
            bookmarks=true,        
            a4paper=true,
            pdftitle={Written with GAPDoc},
            pdfcreator={LaTeX with hyperref package / GAPDoc},
            colorlinks=true,
            backref=page,
            breaklinks=true,
            linkcolor=linkColor,
            citecolor=citeColor,
            filecolor=fileColor,
            urlcolor=urlColor,
            pdfpagemode={UseNone}, 
           ]{hyperref}

\newcommand{\maintitlesize}{\fontsize{50}{55}\selectfont}

% write page numbers to a .pnr log file for online help
\newwrite\pagenrlog
\immediate\openout\pagenrlog =\jobname.pnr
\immediate\write\pagenrlog{PAGENRS := [}
\newcommand{\logpage}[1]{\protect\write\pagenrlog{#1, \thepage,}}
%% were never documented, give conflicts with some additional packages

\newcommand{\GAP}{\textsf{GAP}}

%% nicer description environments, allows long labels
\usepackage{enumitem}
\setdescription{style=nextline}

%% depth of toc
\setcounter{tocdepth}{1}





%% command for ColorPrompt style examples
\newcommand{\gapprompt}[1]{\color{promptColor}{\bfseries #1}}
\newcommand{\gapbrkprompt}[1]{\color{brkpromptColor}{\bfseries #1}}
\newcommand{\gapinput}[1]{\color{gapinputColor}{#1}}


\begin{document}

\logpage{[ 0, 0, 0 ]}
\begin{titlepage}
\mbox{}\vfill

\begin{center}{\maintitlesize \textbf{ IdPGroup \\
\mbox{}}}\\
\vfill

\hypersetup{pdftitle={ IdPGroup }}
\markright{\scriptsize \mbox{}\hfill  IdPGroup  \hfill\mbox{}}
{\Huge \textbf{ Identification functions for different p\texttt{\symbol{45}}groups \\
\mbox{}}}\\
\vfill

{\Huge  1.0 \mbox{}}\\[1cm]
{ 11 August 2026 \mbox{}}\\[1cm]
\mbox{}\\[2cm]
{\Large \textbf{\strut  Henrik Schanze     \strut\mbox{}}}\\
\hypersetup{pdfauthor={ Henrik Schanze    }}
\end{center}\vfill

\mbox{}\\
{\mbox{}\\
\small \noindent \textbf{ Henrik Schanze    }  Email: \href{mailto://henrik.schanze@tu-braunschweig.de} {\texttt{henrik.schanze@tu\texttt{\symbol{45}}braunschweig.de}}\\
  Homepage: \href{https://www.iaa.tu-bs.de/henschan} {\texttt{https://www.iaa.tu\texttt{\symbol{45}}bs.de/henschan}}\\
  Address: \begin{minipage}[t]{8cm}\noindent
 Universitaetsplatz 2\\
 \end{minipage}
}\\
\end{titlepage}

\newpage\setcounter{page}{2}
\newpage

\def\contentsname{Contents\logpage{[ 0, 0, 1 ]}}

\tableofcontents
\newpage

     
\chapter{\textcolor{Chapter }{Introduction}}\label{Chapter_Introduction}
\logpage{[ 1, 0, 0 ]}
\hyperdef{L}{X7DFB63A97E67C0A1}{}
{
  

 The purpose of the \textsf{IdPGroup} package is to implement identification functions for different $p$\texttt{\symbol{45}}groups, which are implemented in the \textsf{SmallGrps} library, but lack an effective identification routine in it. The first
function identifies the groups of order 512. This is done by calculating a
number of invariants of a given group and comparing them with the values of
all possible candidates via a precalculated search tree. More details can be
found in section \ref{Section_Identification}. 

 Section \ref{Section_Load} introduces the definitions of (strong) siblings and includes functions to load
all such sets of groups of certain order. 

 Functions that were used to calculate these lists of groups and which can be
used to calculate them for different orders can be found in section \ref{Section_Calculate}. 

 Once such sets of groups are determined one might want to examine them further
by checking certain properties and calculating additionaly invariants. In
Section \ref{Section_Table} one can find functions to calculate some selected invarants and properties for
such sets and export them into LaTeX\texttt{\symbol{45}}readable lines of
code. 

 Finally section \ref{Section_Build} describes all funcions that were used to build the search tree in the
descirbed identification routine and which can be used to build other such
trees. }

   
\chapter{\textcolor{Chapter }{Functionality}}\label{Chapter_Functionality}
\logpage{[ 2, 0, 0 ]}
\hyperdef{L}{X87F1120883F5B4D0}{}
{
  

 This chapter lists all functions contained in the \textsf{IdPGroup} package. 
\section{\textcolor{Chapter }{Identifying groups via a search tree}}\label{Section_Identification}
\logpage{[ 2, 1, 0 ]}
\hyperdef{L}{X82BA59FA865D19E5}{}
{
  

 This section briefly describes the functions and variables of this package
associated to the identification of an arbitrary group of order 512 via the
search tree \texttt{ID{\textunderscore}GROUP{\textunderscore}512{\textunderscore}TREE} (\ref{IDuScoreGROUPuScore512uScoreTREE}). When saying 'identify a group $G$' we want to find an \emph{id} such that \texttt{SmallGroup(512, id)} is isomorphic to $G$. 

 The used tree is implement as a nested record with the following components: 
\begin{itemize}
\item  "fp": a list of integers, each being the encoded fingerprint of the group
resulting from the specific test performed at the current level of the
identification routine. Each integer is gained from appyling \texttt{IdPGroup{\textunderscore}Pack} (\ref{IdPGroupuScorePack}) to the fingerprint. 
\item  "next": a list comprised of the resprective subbranches corresponding to the
fingerprints in "fp". This is either a list of ids or a branch implemented as
a record itself. 
\end{itemize}
 Depending on the level each record also has one of the following two
components: 
\begin{itemize}
\item  "test": a string that indicates which test was used via \texttt{TestOnGroup} (\ref{TestOnGroup}) to create the listed fingerprints. 
\item  "desc": a list of integers that encode a description of the performed test
using \texttt{EvalFpCoc}. 
\end{itemize}
 

 The first levels of \texttt{ID{\textunderscore}GROUP{\textunderscore}512{\textunderscore}TREE} (\ref{IDuScoreGROUPuScore512uScoreTREE}) try to distinguish groups by the following invariants: 
\begin{itemize}
\item  The rank of the group. 
\item  The collected orders of all elements of the group. 
\item  The numbers of conjugacy classes of each length and representative order (the
order of the elements in the class) of the group. 
\item  The isomorphism type of the parent of the group, i.e. the quotient of the
group and the smallest non\texttt{\symbol{45}}trivial element of the
lower\texttt{\symbol{45}}$p$\texttt{\symbol{45}}central series. 
\item  The abelian invarants, i.e. the isomorphism types of the quotients of the
derived series. 
\item  The cycle structure of the permutations induced on all clusters of conjugacy
classes of same length and representative order by the power map $x\mapsto x^3$. 
\end{itemize}
 

 If necessary these invariants are followed by invaraints resulting from tests
evaluated on partitions of $G$. 

 The initial partition is made of the conjugacy classes of $G$ grouped such that conjugacy classes of the same length and same representative
order belong to the same partitioning subset. The results of \texttt{EvalFpCoc} applied to a partition give the invaraints used to distinguish groups in these
levels of the search tree. The exact behaviour is determined by the content of
the component "desc". After the evaluation of each test, the partition is
refined according to the result of \texttt{EvalFpCoc}. 

 If the preceding invaraints do not lead to a single group, the following (more
time consuming) invaraints are calculated. 
\begin{itemize}
\item  The isomorphism types of the non\texttt{\symbol{45}}trivial factors formed by
the group and elements of the smallest non\texttt{\symbol{45}}trivial element
of the lower\texttt{\symbol{45}}p\texttt{\symbol{45}}central series. 
\item  The isomorphism types of the non\texttt{\symbol{45}}trivial factors formed by
the group and elements of the center. 
\item  The isomorphism types of all maximal subgroups. 
\end{itemize}
 

 Remaining are 139 pairs and 1 triplet of groups of unidentified groups. All
these groups are distinguished by a random isomorphism test. That is done by
selecting random elements from predefined clusters of conjugacy classes, that
were used and potentially refined before, until these elements generate the
group. This list of generators is then turned into a \texttt{pcgs} and the presentation encoded into an integer by the funtion \texttt{CodeGenerators}. The resulting integer modulo 99661 is then compared with precalculated
fingerprints of the candidate groups. If none of the fingerprints coincides
the process is repeated. Thus the runtime of the identification for these
groups is not constant. 

\subsection{\textcolor{Chapter }{IdGroup512}}
\logpage{[ 2, 1, 1 ]}\nobreak
\hyperdef{L}{X81731B9A798C5985}{}
{\noindent\textcolor{FuncColor}{$\triangleright$\enspace\texttt{IdGroup512({\mdseries\slshape G[, opt]})\index{IdGroup512@\texttt{IdGroup512}}
\label{IdGroup512}
}\hfill{\scriptsize (function)}}\\
\textbf{\indent Returns:\ }
The id of the group of order 512 that is isomorphic to \mbox{\texttt{\mdseries\slshape G}}. 



 Identifies the group \mbox{\texttt{\mdseries\slshape G}} of order 512. This function accepts a record \mbox{\texttt{\mdseries\slshape opt}} of options with the following content: 
\begin{itemize}
\item  \mbox{\texttt{\mdseries\slshape tree}}: A search tree that was built using the invariants resulting from $tests$. 
\item  \mbox{\texttt{\mdseries\slshape anupq}}: A boolean that determines if $IsPqIsomorphicPGroup$ from the \textsf{Anupq}\texttt{\symbol{45}}Package should be used to find the to \mbox{\texttt{\mdseries\slshape G}} isomorphic group among the remaining candidates. 
\item  \mbox{\texttt{\mdseries\slshape path}}: A boolean that determines if the path through the tree shall be recorded
while identification and returned as well. 
\end{itemize}
 The default values of \mbox{\texttt{\mdseries\slshape opt}} is \texttt{rec(anupq}:=\texttt{false}, \texttt{tree:=ID{\textunderscore}GROUP{\textunderscore}512{\textunderscore}TREE,
path:=} \texttt{false} \texttt{)}. 
\begin{Verbatim}[commandchars=!@|,fontsize=\small,frame=single,label=Example]
  !gapprompt@gap>| !gapinput@G := Image(IsomorphismPermGroup(SmallGroup(512, 1984)));;|
  !gapprompt@gap>| !gapinput@IdGroup512(G);|
  1984
  !gapprompt@gap>| !gapinput@IdGroup512(G, rec(path := true));|
  [ 1984, rec( fp := [ 2, 8217, 5321 ], pos := [ 2, 188, 1 ],
          tests := [ "Rank", "ElmOrds", [ 40511009 ] ] ) ]
  !gapprompt@gap>| !gapinput@H := Image(IsomorphismPermGroup(SmallGroup(512, 449)));;|
  !gapprompt@gap>| !gapinput@IdGroup512(H, rec(anupq := false));|
  449
  !gapprompt@gap>| !gapinput@IdGroup512(H, rec(anupq := true));|
  449
\end{Verbatim}
 }

 

\subsection{\textcolor{Chapter }{ID{\textunderscore}GROUP{\textunderscore}512{\textunderscore}TREE}}
\logpage{[ 2, 1, 2 ]}\nobreak
\hyperdef{L}{X7E6424F37C14208E}{}
{\noindent\textcolor{FuncColor}{$\triangleright$\enspace\texttt{ID{\textunderscore}GROUP{\textunderscore}512{\textunderscore}TREE\index{ID{\textunderscore}GROUP{\textunderscore}512{\textunderscore}TREE@\texttt{ID{\textunderscore}}\-\texttt{G}\-\texttt{R}\-\texttt{O}\-\texttt{U}\-\texttt{P{\textunderscore}512{\textunderscore}}\-\texttt{TREE}}
\label{IDuScoreGROUPuScore512uScoreTREE}
}\hfill{\scriptsize (global variable)}}\\


 The search tree of groups of order 512 that can determine the vast majority of
groups. At the beginning only a small part is loaded. Other branches are
loaded if necessary. }

 

\subsection{\textcolor{Chapter }{LoadTotalTree}}
\logpage{[ 2, 1, 3 ]}\nobreak
\hyperdef{L}{X7BA3357E86269C3B}{}
{\noindent\textcolor{FuncColor}{$\triangleright$\enspace\texttt{LoadTotalTree({\mdseries\slshape arg})\index{LoadTotalTree@\texttt{LoadTotalTree}}
\label{LoadTotalTree}
}\hfill{\scriptsize (function)}}\\


 Loads the entire search tree \texttt{ID{\textunderscore}GROUP{\textunderscore}512{\textunderscore}TREE} (\ref{IDuScoreGROUPuScore512uScoreTREE}). }

 

\subsection{\textcolor{Chapter }{IdGroup512Exceptions}}
\logpage{[ 2, 1, 4 ]}\nobreak
\hyperdef{L}{X78FA5E5D78E8E44A}{}
{\noindent\textcolor{FuncColor}{$\triangleright$\enspace\texttt{IdGroup512Exceptions({\mdseries\slshape arg})\index{IdGroup512Exceptions@\texttt{IdGroup512Exceptions}}
\label{IdGroup512Exceptions}
}\hfill{\scriptsize (function)}}\\
\textbf{\indent Returns:\ }
A list of bins of group\texttt{\symbol{45}}ids, that \texttt{ID{\textunderscore}GROUP{\textunderscore}512{\textunderscore}TREE} (\ref{IDuScoreGROUPuScore512uScoreTREE}) can not distinguish without random isomorphism tests. 



 
\begin{Verbatim}[commandchars=!@|,fontsize=\small,frame=single,label=Example]
  !gapprompt@gap>| !gapinput@L := IdGroup512Exceptions();;|
  !gapprompt@gap>| !gapinput@Collected(List(L, Size));|
  [ [ 2, 139 ], [ 3, 1 ] ]
  !gapprompt@gap>| !gapinput@L[138];|
  [ 6249609, 6249610, 6249611 ]
\end{Verbatim}
 }

 }

 
\section{\textcolor{Chapter }{Loading (strong) siblings}}\label{Section_Load}
\logpage{[ 2, 2, 0 ]}
\hyperdef{L}{X7EDB790F85B7753E}{}
{
  

 This section describes the functions used to load precalculated lists of
(strong) siblings, i.e. groups that fulfil the following definitions: 
\begin{itemize}
\item  \emph{subgroup\texttt{\symbol{45}}equivalent groups} are groups with a bijection between the sets of proper subgroups s.t. preimage
and image are isomorphic, 
\item  \emph{factor\texttt{\symbol{45}}equivalent groups} are groups with a bijection between the sets of proper factors s.t. preimage
and image are isomorphic, 
\item  \emph{siblings} are subgroup\texttt{\symbol{45}} and factor\texttt{\symbol{45}}equivalent
groups, where the bijection of proper subgroups preservs conjugacy and maps
the Frattini subgroup, lower and upper central series onto the respective
counterparts, 
\item  \emph{strong siblings} are siblings that are also Brauer Pairs. 
\end{itemize}
 All (strong) siblings of orders dividing 512 together with some trivial cases
are implemented in this package and are callable by the following functions. 

\subsection{\textcolor{Chapter }{SiblingsIds}}
\logpage{[ 2, 2, 1 ]}\nobreak
\hyperdef{L}{X7F39858187F267DA}{}
{\noindent\textcolor{FuncColor}{$\triangleright$\enspace\texttt{SiblingsIds({\mdseries\slshape order[, indices]})\index{SiblingsIds@\texttt{SiblingsIds}}
\label{SiblingsIds}
}\hfill{\scriptsize (function)}}\\
\textbf{\indent Returns:\ }
A list of lists of ids of siblings of size \mbox{\texttt{\mdseries\slshape order}}. 



 Loads and returns a list of ids of siblings of size \mbox{\texttt{\mdseries\slshape order}} if available. Currently available sizes are: 
\begin{itemize}
\item  all orders of the form $p\,q$ for $p$ and $q$ primes, 
\item  all powers of two smaller than $2^9$ 
\item  all powers of three smaller than $3^5$ 
\item  all powers of a prime smaller than $p^4 $ 
\item  all products of two powers and 
\item  all orders smaller than 255 except 60, 120, 160, 168, 180, 189, 192 and 240. 
\end{itemize}
 If no list is available for the provided order \texttt{fail} is returned. If a list \mbox{\texttt{\mdseries\slshape indices}} is given, only the respective entries are returned. 
\begin{Verbatim}[commandchars=!@|,fontsize=\small,frame=single,label=Example]
  !gapprompt@gap>| !gapinput@SiblingsIds(101);|
  [ ]
  !gapprompt@gap>| !gapinput@SiblingsIds(128);|
  [ [ 1317, 1322 ], [ 1327, 1329 ], [ 1597, 1598 ] ]
  !gapprompt@gap>| !gapinput@SiblingsIds(256, [1..3]);|
  [ [ 227, 228 ], [ 258, 260, 261 ], [ 845, 846 ] ]
\end{Verbatim}
 }

 

\subsection{\textcolor{Chapter }{Siblings}}
\logpage{[ 2, 2, 2 ]}\nobreak
\hyperdef{L}{X878A3FFE837F9EF6}{}
{\noindent\textcolor{FuncColor}{$\triangleright$\enspace\texttt{Siblings({\mdseries\slshape order[, indices]})\index{Siblings@\texttt{Siblings}}
\label{Siblings}
}\hfill{\scriptsize (function)}}\\
\textbf{\indent Returns:\ }
A list of siblings of size \mbox{\texttt{\mdseries\slshape order}}. 



 Loads and returns a list of siblings of size \mbox{\texttt{\mdseries\slshape order}} if available. If no list is available for the provided order \texttt{fail} is returned. If \mbox{\texttt{\mdseries\slshape indices}} are given, only the respective entries are returned. }

 

\subsection{\textcolor{Chapter }{StrongSiblingsIds}}
\logpage{[ 2, 2, 3 ]}\nobreak
\hyperdef{L}{X8766ECFD817FB4A7}{}
{\noindent\textcolor{FuncColor}{$\triangleright$\enspace\texttt{StrongSiblingsIds({\mdseries\slshape order[, indices]})\index{StrongSiblingsIds@\texttt{StrongSiblingsIds}}
\label{StrongSiblingsIds}
}\hfill{\scriptsize (function)}}\\
\textbf{\indent Returns:\ }
A list of lists of ids of strong siblings of \mbox{\texttt{\mdseries\slshape order}}. 



 Loads and returns a list of ids of strong siblings of \mbox{\texttt{\mdseries\slshape order}} if available. Currently available sizes are: 
\begin{itemize}
\item  all orders of the form $p\,q$ for $p$ and $q$ primes, 
\item  all powers of two smaller than $2^9$ 
\item  all powers of three smaller than $3^5$ 
\item  all powers of a prime smaller than $p^4 $ 
\item  all products of two powers and 
\item  all orders smaller than 255 except 60, 120, 160, 168, 180, 189, 192 and 240. 
\end{itemize}
 If no list is available for the provided order \texttt{fail} is returned. If a list \mbox{\texttt{\mdseries\slshape indices}} is given, only the respective entries are returned. 
\begin{Verbatim}[commandchars=!@|,fontsize=\small,frame=single,label=Example]
  !gapprompt@gap>| !gapinput@StrongSiblingsIds(256);|
  [ [ 1734, 1735 ], [ 1736, 1737 ], [ 1739, 1740 ], [ 1741, 1742 ],
    [ 3678, 3679 ], [ 4154, 4157 ], [ 4155, 4158 ], [ 4156, 4159 ] ]
  !gapprompt@gap>| !gapinput@StrongSiblingsIds(256, [1..3]);|
  [ [ 1734, 1735 ], [ 1736, 1737 ], [ 1739, 1740 ] ]
\end{Verbatim}
 }

 

\subsection{\textcolor{Chapter }{StrongSiblings}}
\logpage{[ 2, 2, 4 ]}\nobreak
\hyperdef{L}{X843063807B06E65B}{}
{\noindent\textcolor{FuncColor}{$\triangleright$\enspace\texttt{StrongSiblings({\mdseries\slshape order[, indices]})\index{StrongSiblings@\texttt{StrongSiblings}}
\label{StrongSiblings}
}\hfill{\scriptsize (function)}}\\
\textbf{\indent Returns:\ }
A list of strong siblings of \mbox{\texttt{\mdseries\slshape order}}. 



 Loads and returns a list of strong siblings of \mbox{\texttt{\mdseries\slshape order}} if available. If no list is available for the provided order \texttt{fail} is returned. If a list \mbox{\texttt{\mdseries\slshape indices}} is given, only the respective entries are returned. }

 }

 
\section{\textcolor{Chapter }{Functions for calculating (strong) siblings}}\label{Section_Calculate}
\logpage{[ 2, 3, 0 ]}
\hyperdef{L}{X7AA188507AFCEC21}{}
{
  

 This section describes all functions for calculating (strong) siblings of a
given order. 

\subsection{\textcolor{Chapter }{SplitBinsViaTest}}
\logpage{[ 2, 3, 1 ]}\nobreak
\hyperdef{L}{X7F0D9E5B7D1947EF}{}
{\noindent\textcolor{FuncColor}{$\triangleright$\enspace\texttt{SplitBinsViaTest({\mdseries\slshape Bins, order, test})\index{SplitBinsViaTest@\texttt{SplitBinsViaTest}}
\label{SplitBinsViaTest}
}\hfill{\scriptsize (function)}}\\
\textbf{\indent Returns:\ }
A list of sublists of the provided lists in \mbox{\texttt{\mdseries\slshape Bins}}. 



 Takes a list of bins with ids of groups of size \mbox{\texttt{\mdseries\slshape order}} and splits them via the outcome of \mbox{\texttt{\mdseries\slshape test}} implemented in \texttt{TestOnGroup}. 
\begin{Verbatim}[commandchars=!@|,fontsize=\small,frame=single,label=Example]
  !gapprompt@gap>| !gapinput@SplitBinsViaTest([[1..10]], 512, "MaxSubId");|
  [ [ 1 ], [ 2 ], [ 3, 8 ], [ 4 ], [ 5 ], [ 6, 7 ], [ 9 ], [ 10 ] ]
\end{Verbatim}
 }

 

\subsection{\textcolor{Chapter }{CalculateSiblings}}
\logpage{[ 2, 3, 2 ]}\nobreak
\hyperdef{L}{X7A81D24885DD46A1}{}
{\noindent\textcolor{FuncColor}{$\triangleright$\enspace\texttt{CalculateSiblings({\mdseries\slshape order[, verbose]})\index{CalculateSiblings@\texttt{CalculateSiblings}}
\label{CalculateSiblings}
}\hfill{\scriptsize (function)}}\\
\textbf{\indent Returns:\ }
A list of lists of ids of groups of size \mbox{\texttt{\mdseries\slshape order}} that are siblings. 



 Determines all siblings among the groups of size \mbox{\texttt{\mdseries\slshape order}} by applying \texttt{SplitBinsViaTest} (\ref{SplitBinsViaTest}) with the tests \texttt{"MaxSubId"}, \texttt{"SpSubId"}, \texttt{"SubId"}, \texttt{"MaxFactId"} and \texttt{"FactId"} iterativly. The tests \texttt{"MaxSubId"} and \texttt{"MaxFactId"} are not necessary, but are efficient and eliminate some sets before applying
the more costly tests. The boolean \mbox{\texttt{\mdseries\slshape verbose}} controls if a collection of bin sizes is printed after each test (\texttt{true} as default). 
\begin{Verbatim}[commandchars=!@|,fontsize=\small,frame=single,label=Example]
  !gapprompt@gap>| !gapinput@CalculateSiblings(128);|
  Test 'MaxSubId' found following bins of size greater one:[ [ 2, 23 ], [ 3, 2 ] ]
  Test 'SpSubId' found following bins of size greater one:[ [ 2, 22 ] ]
  Test 'SubId' found following bins of size greater one:[ [ 2, 16 ] ]
  Test 'MaxFactId' found following bins of size greater one:[ [ 2, 3 ] ]
  Test 'FactId' found following bins of size greater one:[ [ 2, 3 ] ]
  [ [ 1317, 1322 ], [ 1327, 1329 ], [ 1597, 1598 ] ]
  !gapprompt@gap>| !gapinput@CalculateSiblings(128, false);|
  [ [ 1317, 1322 ], [ 1327, 1329 ], [ 1597, 1598 ] ]
\end{Verbatim}
 }

 

\subsection{\textcolor{Chapter }{DetermineBrauerPairs}}
\logpage{[ 2, 3, 3 ]}\nobreak
\hyperdef{L}{X7F8D187184411526}{}
{\noindent\textcolor{FuncColor}{$\triangleright$\enspace\texttt{DetermineBrauerPairs({\mdseries\slshape Bins, order})\index{DetermineBrauerPairs@\texttt{DetermineBrauerPairs}}
\label{DetermineBrauerPairs}
}\hfill{\scriptsize (function)}}\\
\textbf{\indent Returns:\ }
A list of sublists of the provided lists in \mbox{\texttt{\mdseries\slshape Bins}}. 



 Determines the Brauer pairs among the bins of ids of groups of size \mbox{\texttt{\mdseries\slshape order}} listed in \mbox{\texttt{\mdseries\slshape Bins}}. }

 

\subsection{\textcolor{Chapter }{CalculateStrongSiblings}}
\logpage{[ 2, 3, 4 ]}\nobreak
\hyperdef{L}{X7E2FC2447FF1E8FE}{}
{\noindent\textcolor{FuncColor}{$\triangleright$\enspace\texttt{CalculateStrongSiblings({\mdseries\slshape order[, verbose]})\index{CalculateStrongSiblings@\texttt{CalculateStrongSiblings}}
\label{CalculateStrongSiblings}
}\hfill{\scriptsize (function)}}\\
\textbf{\indent Returns:\ }
A list of lists of ids of groups of size \mbox{\texttt{\mdseries\slshape order}} that are strong siblings. 



 First calculats all siblings of size \mbox{\texttt{\mdseries\slshape order}} via \texttt{CalculateSiblings} (\ref{CalculateSiblings}) and then determines all strong siblings among these by \texttt{DetermineBrauerPairs} (\ref{DetermineBrauerPairs}). The boolean \mbox{\texttt{\mdseries\slshape verbose}} controls if a collection of bin sizes is printed after each test while
siblings are calculated (\texttt{true} as default). 
\begin{Verbatim}[commandchars=!@|,fontsize=\small,frame=single,label=Example]
  !gapprompt@gap>| !gapinput@CalculateStrongSiblings(256, false);|
  [ [ 1734, 1735 ], [ 1736, 1737 ], [ 1739, 1740 ], [ 1741, 1742 ], [ 3678, 3679 ],
    [ 4154, 4157 ], [ 4155, 4158 ], [ 4156, 4159 ] ]
\end{Verbatim}
 }

 }

 
\section{\textcolor{Chapter }{Table of bins of groups}}\label{Section_Table}
\logpage{[ 2, 4, 0 ]}
\hyperdef{L}{X792030A679DAAED6}{}
{
  

 This section includes the functions of the package that create and fill a
table for bins of groups, listing certain invaraints and properties and
exporting this table as rows of a latex\texttt{\symbol{45}}tables. 

\subsection{\textcolor{Chapter }{TableOfGroupInvariants}}
\logpage{[ 2, 4, 1 ]}\nobreak
\hyperdef{L}{X7E79E6347A6F8533}{}
{\noindent\textcolor{FuncColor}{$\triangleright$\enspace\texttt{TableOfGroupInvariants({\mdseries\slshape Bins, order})\index{TableOfGroupInvariants@\texttt{TableOfGroupInvariants}}
\label{TableOfGroupInvariants}
}\hfill{\scriptsize (function)}}\\
\textbf{\indent Returns:\ }
A list of records. 



 Calculates some invariants and properties for the groups of order \mbox{\texttt{\mdseries\slshape order}} with ids in lists of \mbox{\texttt{\mdseries\slshape bins}}. If the invariants are equal for all groups of a bin, a record listing all
information is added to the returned list. The calculated invariants are: 
\begin{itemize}
\item  the rank of the groups, 
\item  the $p$\texttt{\symbol{45}}class of the groups, 
\item  the size of the groups of automorphisms. 
\end{itemize}
 Furthermore the following properties are examined for the groups of a bin: 
\begin{itemize}
\item  Are the groups
outer\texttt{\symbol{45}}automorphism\texttt{\symbol{45}}equivaleant, i.e. are
the groups of outer\texttt{\symbol{45}}automorphism isomorphic? 
\item  Are the groups isoclinic? 
\item  Are the groups a modular isomorphic pair (MIP), i.e. are their group rings
over a field of characteristic p isomorphic? 
\end{itemize}
 The last two properties are not examined and left as an empty string in this
function, but can be filled via \texttt{DetermineMIPInTable} (\ref{DetermineMIPInTable}) and \texttt{DetermineIsoclinicInTable} (\ref{DetermineIsoclinicInTable}). 
\begin{Verbatim}[commandchars=!@|,fontsize=\small,frame=single,label=Example]
  !gapprompt@gap>| !gapinput@Bins := StrongSiblingsIds(512, [1, 5, 14]);|
  [ [ 1364, 1368 ], [ 1680, 1681 ], [ 15826, 15827 ] ]
  !gapprompt@gap>| !gapinput@TableOfGroupInvariants(Bins, 512);|
  [ rec( OutEqv := "yes", SizeAut := 8192, ids := [ 1364, 1368 ],
         isoclinic := "",  mip := "", pclass := 5, rank := 2 ),
   rec( OutEqv := "yes", SizeAut := 16384, ids := [ 1680, 1681 ],
        isoclinic := "", mip := "", pclass := 5, rank := 2 ),
   rec( OutEqv := "no", SizeAut := 65536, ids := [ 15826, 15827 ],
        isoclinic := "", mip := "", pclass := 3, rank := 3 ) ]
\end{Verbatim}
 }

 

\subsection{\textcolor{Chapter }{DetermineMIPInTable}}
\logpage{[ 2, 4, 2 ]}\nobreak
\hyperdef{L}{X78D18CBC7C0D8F8C}{}
{\noindent\textcolor{FuncColor}{$\triangleright$\enspace\texttt{DetermineMIPInTable({\mdseries\slshape tab})\index{DetermineMIPInTable@\texttt{DetermineMIPInTable}}
\label{DetermineMIPInTable}
}\hfill{\scriptsize (function)}}\\
\textbf{\indent Returns:\ }
A list of records. 



 Determines if the sets of groups represented by the records in \mbox{\texttt{\mdseries\slshape tab}} are modular isomorphic pairs using the \textsf{ModIsom} package. This information is added to each record of \mbox{\texttt{\mdseries\slshape tab}} and the updated list is returned. This function is included in an extension of
the package and only available if the respective package is loaded. }

 

\subsection{\textcolor{Chapter }{DetermineIsoclinicInTable}}
\logpage{[ 2, 4, 3 ]}\nobreak
\hyperdef{L}{X7A5F26D17E268E4B}{}
{\noindent\textcolor{FuncColor}{$\triangleright$\enspace\texttt{DetermineIsoclinicInTable({\mdseries\slshape tab})\index{DetermineIsoclinicInTable@\texttt{DetermineIsoclinicInTable}}
\label{DetermineIsoclinicInTable}
}\hfill{\scriptsize (function)}}\\
\textbf{\indent Returns:\ }
A list of records. 



 Determines if the sets of groups represented by the records in \mbox{\texttt{\mdseries\slshape tab}} are isoclinic using the \textsf{XMod} package. This information is added to each record and the of \mbox{\texttt{\mdseries\slshape tab}} updated list is returned. This function is included in an extension of the
package and only available if the respective package is loaded. }

 

\subsection{\textcolor{Chapter }{MergeTablesOfInvariants}}
\logpage{[ 2, 4, 4 ]}\nobreak
\hyperdef{L}{X83C7850B7F8202AA}{}
{\noindent\textcolor{FuncColor}{$\triangleright$\enspace\texttt{MergeTablesOfInvariants({\mdseries\slshape tab1, tab2})\index{MergeTablesOfInvariants@\texttt{MergeTablesOfInvariants}}
\label{MergeTablesOfInvariants}
}\hfill{\scriptsize (function)}}\\
\textbf{\indent Returns:\ }
A list of records. 



 Merges the tables \mbox{\texttt{\mdseries\slshape tab1}} and \mbox{\texttt{\mdseries\slshape tab2}} into one table. If a record entry of only one is the empty string the
non\texttt{\symbol{45}}empty string is used. This enables the parallel
computation of e.g. \texttt{DetermineMIPInTable} (\ref{DetermineMIPInTable}) and \texttt{DetermineIsoclinicInTable} (\ref{DetermineIsoclinicInTable}). }

 

\subsection{\textcolor{Chapter }{PrintTableToFile}}
\logpage{[ 2, 4, 5 ]}\nobreak
\hyperdef{L}{X79CD27767C5C091C}{}
{\noindent\textcolor{FuncColor}{$\triangleright$\enspace\texttt{PrintTableToFile({\mdseries\slshape tab, filename})\index{PrintTableToFile@\texttt{PrintTableToFile}}
\label{PrintTableToFile}
}\hfill{\scriptsize (function)}}\\


 Exports the content of the table \mbox{\texttt{\mdseries\slshape tab}} into "\mbox{\texttt{\mdseries\slshape filename}}{\textunderscore}table.txt" as rows of a latex table. }

 }

 
\section{\textcolor{Chapter }{Functions for building a search tree}}\label{Section_Build}
\logpage{[ 2, 5, 0 ]}
\hyperdef{L}{X7E8276BF79FDD898}{}
{
  

 This section describes all functions and variables connected to building a
search tree that can be used with \texttt{IdGroup512} (\ref{IdGroup512}) to identify groups via implemented invariants. 

 In short this is done be initializing a tree via calling \texttt{TestOnClusterList} (\ref{TestOnClusterList}) with a selected test on a list of (lists of) ids of groups of given order and
then iteratnig the following three steps 
\begin{itemize}
\item  Collecting the clusters of the curent tree via \texttt{ClusterList} (\ref{ClusterList}), 
\item  Performing a choosen test on these clusters via \texttt{TestOnClusterList} (\ref{TestOnClusterList}), 
\item  Extending the tree with these test results via \texttt{AddClusterTestResults} (\ref{AddClusterTestResults}). 
\end{itemize}
 A very effective set of invaraints can be obtained by calling \texttt{CocTestOnClusterList} (\ref{CocTestOnClusterList}) in the second step. This should only be done when the groups of each provided
cluster yields the same value when the test \texttt{CocList} is applied. Furthermore these invaraints should be used for as few as possible
clusters since \texttt{CocTestOnClusterList} (\ref{CocTestOnClusterList}) includes a search for suitable parameters than can take quite some time. 

\subsection{\textcolor{Chapter }{IdPGroup{\textunderscore}Pack}}
\logpage{[ 2, 5, 1 ]}\nobreak
\hyperdef{L}{X7894781585B18100}{}
{\noindent\textcolor{FuncColor}{$\triangleright$\enspace\texttt{IdPGroup{\textunderscore}Pack({\mdseries\slshape list})\index{IdPGroup{\textunderscore}Pack@\texttt{IdPGroup{\textunderscore}Pack}}
\label{IdPGroupuScorePack}
}\hfill{\scriptsize (function)}}\\
\textbf{\indent Returns:\ }
An interger encoding the content of \mbox{\texttt{\mdseries\slshape list}}. 



 Encodes the content of the provided \mbox{\texttt{\mdseries\slshape list}} into an integer. 
\begin{Verbatim}[commandchars=!@|,fontsize=\small,frame=single,label=Example]
  !gapprompt@gap>| !gapinput@IdPGroup_Pack([1,2,3]);|
  123
  !gapprompt@gap>| !gapinput@IdPGroup_Pack([606, 109, 1405]);|
  63095
  !gapprompt@gap>| !gapinput@IdPGroup_Pack([[1,2,3], [606, 1405]]);|
  19765
\end{Verbatim}
 }

 

\subsection{\textcolor{Chapter }{ClusterConjClasses}}
\logpage{[ 2, 5, 2 ]}\nobreak
\hyperdef{L}{X828C60087DC9E1BD}{}
{\noindent\textcolor{FuncColor}{$\triangleright$\enspace\texttt{ClusterConjClasses({\mdseries\slshape G})\index{ClusterConjClasses@\texttt{ClusterConjClasses}}
\label{ClusterConjClasses}
}\hfill{\scriptsize (function)}}\\
\textbf{\indent Returns:\ }
A list of lists containing the conjugacy classes of \mbox{\texttt{\mdseries\slshape G}}. 



 Computes the conjugacy classes of the group \mbox{\texttt{\mdseries\slshape G}} and clusters them such that each clusters contains all classes of a certain
length and certain representative order (meaning the order of any
representative). }

 

\subsection{\textcolor{Chapter }{PowerMapCocCycSt}}
\logpage{[ 2, 5, 3 ]}\nobreak
\hyperdef{L}{X7D29A07883DABFD7}{}
{\noindent\textcolor{FuncColor}{$\triangleright$\enspace\texttt{PowerMapCocCycSt({\mdseries\slshape coc, n})\index{PowerMapCocCycSt@\texttt{PowerMapCocCycSt}}
\label{PowerMapCocCycSt}
}\hfill{\scriptsize (function)}}\\
\textbf{\indent Returns:\ }
A lists of lists of lists of integers 



 Examines the permutation induced by the power map $x\mapsto x^n$ on each cluster of \mbox{\texttt{\mdseries\slshape coc}}. If a class of a cluster is fixed, the permutation within the class is
determined and added to the lists of permutations of the cluster. Finally the
cycle structure of each permutation associated to a cluster is examined via \texttt{CycleStructurePerm} and saved as a list of integers (undound entries are replaced by 0). }

 

\subsection{\textcolor{Chapter }{ClusterList}}
\logpage{[ 2, 5, 4 ]}\nobreak
\hyperdef{L}{X7812E18678076D7B}{}
{\noindent\textcolor{FuncColor}{$\triangleright$\enspace\texttt{ClusterList({\mdseries\slshape tree})\index{ClusterList@\texttt{ClusterList}}
\label{ClusterList}
}\hfill{\scriptsize (function)}}\\
\textbf{\indent Returns:\ }
A list of all clusters of groups in \mbox{\texttt{\mdseries\slshape tree}} (i.e. the leafs of \mbox{\texttt{\mdseries\slshape tree}}). 



 Iterates through \mbox{\texttt{\mdseries\slshape tree}} and collects information about each leaf as a record. This includes the size
of the leaf, the path through \mbox{\texttt{\mdseries\slshape tree}} leading to the leaf, the tests performed including descriptions of performed \texttt{EvalFpCoc}\texttt{\symbol{45}}tests and the ids of the groups contained in the leaf. 
\begin{Verbatim}[commandchars=!@|,fontsize=\small,frame=single,label=Example]
  !gapprompt@gap>| !gapinput@tree := rec( fp := [ 63049, 75721, 4310 ],  next := [ [ 3, 8 ], [ 4 ],|
              rec( fp := [ 11744, 86646 ], next := [ [ 5 ], [ 6, 7 ] ],
                  test := "MaxSubId" ) ], test := "ElmOrds" );;
  !gapprompt@gap>| !gapinput@ClusterList(tree);|
  [ rec( descl := [  ], ids := [ 3, 8 ], path := [ 63049 ], size := 2,
         tests := [ "ElmOrds" ] ),
    rec( descl := [  ], ids := [ 4 ], path := [ 75721 ], size := 1,
         tests := [ "ElmOrds" ] ),
    rec( descl := [  ], ids := [ 5 ], path := [ 4310, 11744 ], size := 1,
         tests := [ "ElmOrds", "MaxSubId" ] ),
    rec( descl := [  ], ids := [ 6, 7 ], path := [ 4310, 86646 ], size := 2,
         tests := [ "ElmOrds", "MaxSubId" ] ) ]
\end{Verbatim}
 }

 

\subsection{\textcolor{Chapter }{TestOnGroup}}
\logpage{[ 2, 5, 5 ]}\nobreak
\hyperdef{L}{X7C14B8DF7A4C0115}{}
{\noindent\textcolor{FuncColor}{$\triangleright$\enspace\texttt{TestOnGroup({\mdseries\slshape G, test[, opt]})\index{TestOnGroup@\texttt{TestOnGroup}}
\label{TestOnGroup}
}\hfill{\scriptsize (function)}}\\
\textbf{\indent Returns:\ }
The result of \mbox{\texttt{\mdseries\slshape test}} applied to \mbox{\texttt{\mdseries\slshape G}}, if desried encoded as an integer. 



 Performs the selected \mbox{\texttt{\mdseries\slshape test}} on the group \mbox{\texttt{\mdseries\slshape G}}. The following choices for \mbox{\texttt{\mdseries\slshape test}} are implemented: 
\begin{itemize}
\item  \mbox{\texttt{\mdseries\slshape "ElmOrds"}}: Collects the orders of all elements of \mbox{\texttt{\mdseries\slshape G}}. 
\item  \mbox{\texttt{\mdseries\slshape "Center"}}: Calculates the abelian invariants of the center of \mbox{\texttt{\mdseries\slshape G}}. 
\item  \mbox{\texttt{\mdseries\slshape "Rank"}}: Calculates the rank of \mbox{\texttt{\mdseries\slshape G}}. 
\item  \mbox{\texttt{\mdseries\slshape "AbelInv"}}: Computes the order of each term of the derived series and the abelian
Invariants of the derived factors of \mbox{\texttt{\mdseries\slshape G}}. 
\item  \mbox{\texttt{\mdseries\slshape "UCSId"}}: Identifies the proper subgroups of the upper central series of \mbox{\texttt{\mdseries\slshape G}} via \texttt{IdGroup}. 
\item  \mbox{\texttt{\mdseries\slshape "LCSId"}}: Identifies the proper subgroups of the lower central series of \mbox{\texttt{\mdseries\slshape G}} via \texttt{IdGroup}. 
\item  \mbox{\texttt{\mdseries\slshape "PCSId"}}: Identifies the proper subgroups of the p\texttt{\symbol{45}}central series
of \mbox{\texttt{\mdseries\slshape G}} via \texttt{IdGroup}. 
\item  \mbox{\texttt{\mdseries\slshape "SpSubId"}}: Computes the Frattini subgroup, the upper and lower central series of \mbox{\texttt{\mdseries\slshape G}} and identifies all terms if possible via \texttt{IdGroup}. 
\item  \mbox{\texttt{\mdseries\slshape "MaxSubId"}}: Collects the sizes and isomorphism types of all maximal subgroups of \mbox{\texttt{\mdseries\slshape G}} via \texttt{IdGroup}. 
\item  \mbox{\texttt{\mdseries\slshape "SubId"}}: Collects the sizes and isormorphism types of conjugacy classes of all proper
subgoups via \texttt{IdGroup}. 
\item  \mbox{\texttt{\mdseries\slshape "MaxFactId"}}: Collects the sizes and isomorphism types of all maximal factors of \mbox{\texttt{\mdseries\slshape G}} via \texttt{IdGroup}. 
\item  \mbox{\texttt{\mdseries\slshape "CentQuot"}}: Collects the sizes and isomorphism types of all quotients of \mbox{\texttt{\mdseries\slshape G}} with non\texttt{\symbol{45}}trivial elements of the center of \mbox{\texttt{\mdseries\slshape G}} via \texttt{IdGroup}. 
\item  \mbox{\texttt{\mdseries\slshape "LCSFactId"}}: Collects the isomorphism types of all quotients of \mbox{\texttt{\mdseries\slshape G}} with with cyclic subgroups generated by by non\texttt{\symbol{45}}trivial
elements of the last non\texttt{\symbol{45}}trivial term of the
p\texttt{\symbol{45}}central series. 
\item  \mbox{\texttt{\mdseries\slshape "FactId"}}: Identifies all proper factors of \mbox{\texttt{\mdseries\slshape G}} via \texttt{IdGroup}. 
\item  \mbox{\texttt{\mdseries\slshape "CocList"}}: Collects the number of conjugacy classes of a \mbox{\texttt{\mdseries\slshape G}} of same length and same representative order (the order of a representative). 
\item  \mbox{\texttt{\mdseries\slshape "PMCycSt"}}: Determines the cycle\texttt{\symbol{45}}structure invariant for the power
map $x\mapsto x^3$ on the clusters of conjugacy classes via \texttt{PowerMapCocCycSt}. 
\item  \mbox{\texttt{\mdseries\slshape "OutId"}}: Tries to identify the groups of outer automorphisms of \mbox{\texttt{\mdseries\slshape G}}, either via \texttt{IdGroup} (if available) or via its size and \texttt{TestOnGroup} with \mbox{\texttt{\mdseries\slshape "ElmOrds"}} and \mbox{\texttt{\mdseries\slshape "CocList"}} (\mbox{\texttt{\mdseries\slshape "MaxSubId"}} and \mbox{\texttt{\mdseries\slshape "CentQuot"}} can also be used if feasible). 
\item  \mbox{\texttt{\mdseries\slshape "LattEqv"}}: Calculates all conjugacy classes of proper subgoups of \mbox{\texttt{\mdseries\slshape G}}, stores their size, identifies the the isomorphism types of each class and
determines the isomorphism types and conjugacy classes in \mbox{\texttt{\mdseries\slshape G}} of its maximal subgroups and minimal supergroups. 
\end{itemize}
 This function accepts a record \mbox{\texttt{\mdseries\slshape opt}} of options with the single content \mbox{\texttt{\mdseries\slshape opt.pack}}, a boolean that determines if \texttt{IdPGroup{\textunderscore}Pack} should be applied to the result of a test. The default value of \mbox{\texttt{\mdseries\slshape opt}} is $rec(pack:=$\texttt{true}$)$. 
\begin{Verbatim}[commandchars=!@|,fontsize=\small,frame=single,label=Example]
  !gapprompt@gap>| !gapinput@G := SmallGroup(512, 1405);;|
  !gapprompt@gap>| !gapinput@TestOnGroup(G, "ElmOrds");|
  85416
  !gapprompt@gap>| !gapinput@TestOnGroup(G, "ElmOrds", rec(pack := false));|
  [ [ 1, 1 ], [ 2, 7 ], [ 4, 24 ], [ 8, 288 ], [ 16, 64 ], [ 32, 128 ] ]
\end{Verbatim}
 }

 

\subsection{\textcolor{Chapter }{TestOnClusterList}}
\logpage{[ 2, 5, 6 ]}\nobreak
\hyperdef{L}{X7BC4B6397C17B535}{}
{\noindent\textcolor{FuncColor}{$\triangleright$\enspace\texttt{TestOnClusterList({\mdseries\slshape CL, test[, opt]})\index{TestOnClusterList@\texttt{TestOnClusterList}}
\label{TestOnClusterList}
}\hfill{\scriptsize (function)}}\\
\textbf{\indent Returns:\ }
Either a list of pairs each containg the path to the respective cluster (which
is a leaf of a search tree) and the subbranch resulting from the results of \mbox{\texttt{\mdseries\slshape test}}, or just the list of subbranches if \mbox{\texttt{\mdseries\slshape CL}} contains only lists of integers. If \mbox{\texttt{\mdseries\slshape CL}} consits of a single list of integers the single resulting subbranch is
returned on its own. 



 Performs the choosen test on the groups of each record in the \mbox{\texttt{\mdseries\slshape CL}}. The tests are performed by \texttt{TestOnGroup} (\ref{TestOnGroup}). \mbox{\texttt{\mdseries\slshape CL}} could either be a list of records produced by \texttt{ClusterList} (\ref{ClusterList}) or be a list of lists containing ids. In the former case the output is a list
of pairs containg the path of a cluster and the subbranch resulting from the
test. In the latter case only a list of subbranches is returned. If in the
latter case the returned list only one contains one subbranch, the subbranch
itself is returned. This is useful when building the first level of a search
tree. This function accepts a record \mbox{\texttt{\mdseries\slshape opt}} of options with the following content: 
\begin{itemize}
\item  \mbox{\texttt{\mdseries\slshape pack}}: A boolean that is handed to \texttt{TestOnGroup} (\ref{TestOnGroup}). 
\item  \mbox{\texttt{\mdseries\slshape ign1cl}}: A boolean that determines if clusters of size 1 should be ignored. 
\item  \mbox{\texttt{\mdseries\slshape order}}: An integer giving the order of the groups with ids in \mbox{\texttt{\mdseries\slshape CL}}. 
\item  \mbox{\texttt{\mdseries\slshape verbose}}: A boolean that determines if progress should be printed. 
\item  \mbox{\texttt{\mdseries\slshape notetest}}: A boolean that decides if the performed test should be noted in the returned
subbranches. If so subbranches, that do not split a cluster are omited. 
\end{itemize}
 The default value of \mbox{\texttt{\mdseries\slshape opt}} is $rec(pack$:=\texttt{true}, $ign1cl$:=\texttt{true}, $order$:=512, $verbose$:=\texttt{true}, $notetest$:=\texttt{true}$)$. 
\begin{Verbatim}[commandchars=!@|,fontsize=\small,frame=single,label=Example]
  !gapprompt@gap>| !gapinput@TestOnClusterList([[1..10]],"ElmOrds", rec(verbose:=false));|
  rec( fp := [ 88570, 4310, 63049, 75721 ],
       next := [ [ 1 ], [ 2, 5, 6, 7, 9, 10 ], [ 3, 8 ], [ 4 ] ],
       test := "ElmOrds")
  !gapprompt@gap>| !gapinput@TestOnClusterList([[1,2,5],[7,8,9]],"ElmOrds", rec(verbose:=false));|
  [ rec( fp := [ 88570, 4310 ], next := [ [ 1 ], [ 2, 5 ] ],
         test := "ElmOrds" ),
    rec( fp := [ 4310, 63049 ], next := [ [ 7, 9 ], [ 8 ] ],
         test := "ElmOrds") ]
  !gapprompt@gap>| !gapinput@tree := TestOnClusterList([[3..8]],"ElmOrds", rec(verbose:=false));;|
  !gapprompt@gap>| !gapinput@CL := ClusterList(tree);|
  [ rec( descl := [  ], ids := [ 3, 8 ], path := [ 63049 ], size := 2,
         tests := [ "ElmOrds" ]),
    rec( descl := [  ], ids := [ 4 ], path := [ 75721 ], size := 1,
         tests := [ "ElmOrds" ]),
    rec( descl := [  ], ids := [ 5, 6, 7 ], path := [ 4310 ], size := 3,
         tests := [ "ElmOrds" ]) ]
  !gapprompt@gap>| !gapinput@TestOnClusterList(CL, "MaxSubId", rec(verbose:=false, notetest := false))|
  [ [ [ 63049 ], rec( fp := [ 62603 ], next := [ [ 3, 8 ] ] ) ],
    [ [ 4310 ], rec( fp := [ 11744, 86646 ], next := [ [ 5 ], [ 6, 7 ] ] ) ] ]
\end{Verbatim}
 }

 

\subsection{\textcolor{Chapter }{AddClusterTestResults}}
\logpage{[ 2, 5, 7 ]}\nobreak
\hyperdef{L}{X85F0BFE88445D1B7}{}
{\noindent\textcolor{FuncColor}{$\triangleright$\enspace\texttt{AddClusterTestResults({\mdseries\slshape tree, TR})\index{AddClusterTestResults@\texttt{AddClusterTestResults}}
\label{AddClusterTestResults}
}\hfill{\scriptsize (function)}}\\
\textbf{\indent Returns:\ }
The extended \mbox{\texttt{\mdseries\slshape tree}} with the subbranches in \mbox{\texttt{\mdseries\slshape TR}}. 



 Extends \mbox{\texttt{\mdseries\slshape tree}} by the test results stored in \mbox{\texttt{\mdseries\slshape TR}}, a list of pairs each containing the path in \mbox{\texttt{\mdseries\slshape tree}} and a branch (a tree of depth one), extending the corresponding cluster of the
inital \mbox{\texttt{\mdseries\slshape tree}}. 
\begin{Verbatim}[commandchars=!@|,fontsize=\small,frame=single,label=Example]
  !gapprompt@gap>| !gapinput@tree := TestOnClusterList([[3..8]],"ElmOrds", rec(verbose:=false));|
  rec( fp := [ 63049, 75721, 4310 ], next := [ [ 3, 8 ], [ 4 ], [ 5, 6, 7 ] ],
       test := "ElmOrds")
  !gapprompt@gap>| !gapinput@TR:=TestOnClusterList(ClusterList(tree),"MaxSubId",rec(verbose:=false));|
  [ [ [ 4310 ], rec( fp := [ 11744, 86646 ], next := [ [ 5 ], [ 6, 7 ] ],
                     test := "MaxSubId" ) ] ]
  !gapprompt@gap>| !gapinput@AddClusterTestResults(tree,TR);|
  rec( fp := [ 63049, 75721, 4310 ],
       next := [ rec( fp := [ 62603 ], next := [ [ 3, 8 ] ] ), [ 4 ],
                 rec( fp := [ 11744, 86646 ], next := [ [ 5 ], [ 6, 7 ] ],
       test := "MaxSubId") ], test := "ElmOrds" )
\end{Verbatim}
 }

 

\subsection{\textcolor{Chapter }{BuildSearchTree}}
\logpage{[ 2, 5, 8 ]}\nobreak
\hyperdef{L}{X82CDC04E8588770B}{}
{\noindent\textcolor{FuncColor}{$\triangleright$\enspace\texttt{BuildSearchTree({\mdseries\slshape N, order, tests})\index{BuildSearchTree@\texttt{BuildSearchTree}}
\label{BuildSearchTree}
}\hfill{\scriptsize (function)}}\\
\textbf{\indent Returns:\ }
A search tree (implemented as a record). 



 Builds a search tree for the groups with ids in \mbox{\texttt{\mdseries\slshape N}} or, if \mbox{\texttt{\mdseries\slshape N}} is an integer, the first \mbox{\texttt{\mdseries\slshape N}} groups of size \mbox{\texttt{\mdseries\slshape order}}, using \texttt{TestOnClusterList} (\ref{TestOnClusterList}) with the tests from the list \mbox{\texttt{\mdseries\slshape tests}}. 
\begin{Verbatim}[commandchars=!@|,fontsize=\small,frame=single,label=Example]
  !gapprompt@gap>| !gapinput@BuildSearchTree(5,512,["ElmOrds", "MaxSubId"]);|
  rec( fp := [ 88570, 4310, 63049, 75721 ], next := [ [ 1 ],
       rec( fp := [ 76646, 11744 ], next := [ [ 2 ], [ 5 ] ],
       test := "MaxSubId" ), [ 3 ], [ 4 ] ], test := "ElmOrds" )
  !gapprompt@gap>| !gapinput@BuildSearchTree([1,2,3,5,7], 512, ["ElmOrds", "CocList"]);|
  rec( fp := [ 88570, 4310, 63049 ], next := [ [ 1 ],
       rec( fp := [ 96561, 32173 ], next := [ [ 2 ], [ 5, 7 ] ],
       test := "CocList" ), [ 3 ] ], test := "ElmOrds" )
\end{Verbatim}
 }

 

\subsection{\textcolor{Chapter }{FindParametersForCocTest}}
\logpage{[ 2, 5, 9 ]}\nobreak
\hyperdef{L}{X7D2DA3AB79231078}{}
{\noindent\textcolor{FuncColor}{$\triangleright$\enspace\texttt{FindParametersForCocTest({\mdseries\slshape Coc, paralist, it, maxit[, opt]})\index{FindParametersForCocTest@\texttt{FindParametersForCocTest}}
\label{FindParametersForCocTest}
}\hfill{\scriptsize (function)}}\\
\textbf{\indent Returns:\ }
A pair made of a list of return values from \texttt{EvalFpCoc} and a list describing the test performed by \texttt{EvalFpCoc}. 



 Does a parameter search such that evaluating \texttt{EvalFpCoc} on the clusters of conjugacy classes in \mbox{\texttt{\mdseries\slshape coc}} results in at least two different values. If no suitable parameters are found,
the function calls itself iterativly on parameters that alter \mbox{\texttt{\mdseries\slshape Coc}} until either a suitable list of parameters is found or \mbox{\texttt{\mdseries\slshape it}} is equal to \mbox{\texttt{\mdseries\slshape maxit}}. The parameters of the current iteration are provided in the list \mbox{\texttt{\mdseries\slshape paralist}}. If the function is initialized with \mbox{\texttt{\mdseries\slshape it}}=0, it only returns a list of parameters that alter the clusters in \mbox{\texttt{\mdseries\slshape Coc}}. Depending on the size of \mbox{\texttt{\mdseries\slshape coc}} their can be a lot of parameters to check. To decrease the needed time the
following configurations are skipped: 
\begin{itemize}
\item  $[1,i,j]$, where the representative order of \mbox{\texttt{\mdseries\slshape coc[i]}} and j are not coprime. 
\item  $[2,i,j]$, where the representative order of \mbox{\texttt{\mdseries\slshape coc[i]}} is less or equal then the order of \mbox{\texttt{\mdseries\slshape coc[j]}} or the length of the classes in \mbox{\texttt{\mdseries\slshape coc[i]}} is less then the ones of \mbox{\texttt{\mdseries\slshape coc[j]}}. 
\item  $[3,i,j]$, where the representative order of \mbox{\texttt{\mdseries\slshape coc[i]}} is larger or equal then the order of \mbox{\texttt{\mdseries\slshape coc[j]}} or the length of the classes in \mbox{\texttt{\mdseries\slshape coc[i]}} is larger then the ones of \mbox{\texttt{\mdseries\slshape coc[j]}}. 
\item  $[4,i,j,k]$, where either the classes in \mbox{\texttt{\mdseries\slshape coc[i]}} or \mbox{\texttt{\mdseries\slshape coc[j]}} are of length one. 
\item  $[5,i,j,k]$, where either the classes in \mbox{\texttt{\mdseries\slshape coc[i]}} or \mbox{\texttt{\mdseries\slshape coc[j]}} are of length one. 
\end{itemize}
 This function accepts a record \mbox{\texttt{\mdseries\slshape opt}} of options with the following content: 
\begin{itemize}
\item  \mbox{\texttt{\mdseries\slshape target}}: A list of integers that determines the desired second entries of the
parameters. 
\item  \mbox{\texttt{\mdseries\slshape maxlevel}}: An integer that determines the maximal first parameter entry. That might be
useful, since the parameters for test 4 and 5 often need more time then the
rest. 
\end{itemize}
 The default value of \mbox{\texttt{\mdseries\slshape opt}} is $rec(target:=[1..Size(Coc[1])], maxlevel:=5$. 
\begin{Verbatim}[commandchars=!@|,fontsize=\small,frame=single,label=Example]
  !gapprompt@gap>| !gapinput@Coc:=List([3, 8],id->ClusterConjClasses(SmallGroup(512, id)));;|
  !gapprompt@gap>| !gapinput@FindParametersForCocTest(Coc, [], 1, 1);|
  [ [ 1363, 46760 ], [ 3, 2, 4 ] ]
  !gapprompt@gap>| !gapinput@FindParametersForCocTest(Coc, [], 1, 1, rec(target:=[4]));|
  [ [ 63010, 76026 ], [ [ 4, 4, 6, 5 ] ] ]
  !gapprompt@gap>| !gapinput@FindParametersForCocTest(Coc, [], 0, 1, rec(target:=[4]));|
  [ [ 3, 4, 6 ], [ 4, 4, 6, 3 ], [ 4, 4, 6, 4 ], [ 4, 4, 6, 5 ] ]
  !gapprompt@gap>| !gapinput@FindParametersForCocTest(Coc, [], 0, 1, rec(target:=[4], maxlevel := 3));|
  [ [ 3, 4, 6 ] ]
  !gapprompt@gap>| !gapinput@Coc:=List([150, 151],id->ClusterConjClasses(SmallGroup(512, id)));;|
  !gapprompt@gap>| !gapinput@FindParametersForCocTest(Coc, [], 1, 2);|
  [ [ 42, 1626 ], [ [ 1, 7, 5 ], [ 2, 11, 7 ] ] ]
\end{Verbatim}
 }

 

\subsection{\textcolor{Chapter }{DeEncodeDesc}}
\logpage{[ 2, 5, 10 ]}\nobreak
\hyperdef{L}{X833070DB8585CF82}{}
{\noindent\textcolor{FuncColor}{$\triangleright$\enspace\texttt{DeEncodeDesc({\mdseries\slshape desc})\index{DeEncodeDesc@\texttt{DeEncodeDesc}}
\label{DeEncodeDesc}
}\hfill{\scriptsize (function)}}\\
\textbf{\indent Returns:\ }
An interger or a list of integers 



 De\texttt{\symbol{45}} or encodes a description for \texttt{EvalFPCoc}, depending on if an integer (decodes to a list) or a list of integers
(encodes to an integer) is provided. 
\begin{Verbatim}[commandchars=!@|,fontsize=\small,frame=single,label=Example]
  !gapprompt@gap>| !gapinput@DeEncodeDesc([3,2,4]);|
  302004
  !gapprompt@gap>| !gapinput@DeEncodeDesc(302004);|
  [3, 2, 4]
\end{Verbatim}
 }

 

\subsection{\textcolor{Chapter }{AlterCocAccordingToFp}}
\logpage{[ 2, 5, 11 ]}\nobreak
\hyperdef{L}{X85397672793CECF9}{}
{\noindent\textcolor{FuncColor}{$\triangleright$\enspace\texttt{AlterCocAccordingToFp({\mdseries\slshape coc, desc, fp})\index{AlterCocAccordingToFp@\texttt{AlterCocAccordingToFp}}
\label{AlterCocAccordingToFp}
}\hfill{\scriptsize (function)}}\\
\textbf{\indent Returns:\ }
A list of lists containing the conjugacy classes of a group. 



 Alters the clusters of conjugacy classes \mbox{\texttt{\mdseries\slshape coc}} according the result \mbox{\texttt{\mdseries\slshape fp}} from \texttt{EvalFPCoc} evaluated with \mbox{\texttt{\mdseries\slshape desc}} by spliting the cluster at position \mbox{\texttt{\mdseries\slshape desc[2]}} of \mbox{\texttt{\mdseries\slshape coc}} into smaller ones. 
\begin{Verbatim}[commandchars=!@|,fontsize=\small,frame=single,label=Example]
  !gapprompt@gap>| !gapinput@coc := ClusterConjClasses(SmallGroup(512, 3));;|
  !gapprompt@gap>| !gapinput@List(coc, Size);|
  [ 1, 15, 4, 24, 12, 48 ]
  !gapprompt@gap>| !gapinput@fp := EvalFpCoc(coc, [3,2,4]);|
  [ 0, 0, 16, 16, 0, 0, 0, 0, 0, 16, 0, 0, 0, 0, 0 ]
  !gapprompt@gap>| !gapinput@cocnew := AlterCocAccordingToFp(coc,[3,2,4],fp);;|
  !gapprompt@gap>| !gapinput@List(cocnew, Size);|
  [ 1, 12, 3, 4, 24, 12, 48 ]
\end{Verbatim}
 }

 

\subsection{\textcolor{Chapter }{CocSplitRoutine}}
\logpage{[ 2, 5, 12 ]}\nobreak
\hyperdef{L}{X863A18A386DF3E70}{}
{\noindent\textcolor{FuncColor}{$\triangleright$\enspace\texttt{CocSplitRoutine({\mdseries\slshape Ids, order, desclist, iterate})\index{CocSplitRoutine@\texttt{CocSplitRoutine}}
\label{CocSplitRoutine}
}\hfill{\scriptsize (function)}}\\
\textbf{\indent Returns:\ }
A search tree (implemented as a record). 



 Tries to split the list \mbox{\texttt{\mdseries\slshape Ids}} of group ids of \mbox{\texttt{\mdseries\slshape order}} with ids as far as possible on the returned values of \texttt{EvalFpCoc} and returns the resulting subbranch for a search tree with the relevant
details. The parameter search is done with at most \mbox{\texttt{\mdseries\slshape iterate}} iterations. Past perfomances of \texttt{EvalFpCoc} which altered the clusters of conjugacy classes of the groups can be recreated
in advance via the list \mbox{\texttt{\mdseries\slshape desclist}}. If no such splitis possible a warning is printed and \mbox{\texttt{\mdseries\slshape Ids}} is returned. 
\begin{Verbatim}[commandchars=!@|,fontsize=\small,frame=single,label=Example]
  !gapprompt@gap>| !gapinput@CocSplitRoutine([5,6,7], 512, [], 1);|
  rec( desc := [104003], fp := [ 1214, 16 ], next := [ [ 5 ],
       rec( desc := [40204005], fp := [ 20792, 46881 ],
            next := [ [ 6 ], [ 7 ] ] ) ] )
  !gapprompt@gap>| !gapinput@CocSplitRoutine([1529,1530,1531], 512, [], 2);|
  rec( desc := [ 209003 ], fp := [ 4, 14 ], next :=
     [ rec( desc := [ 107003, 212007 ], fp := [ 14, 618 ],
      next := [ [ 1529 ], [ 1530 ] ] ), [ 1531 ] ] )
\end{Verbatim}
 }

 

\subsection{\textcolor{Chapter }{CocTestOnClusterList}}
\logpage{[ 2, 5, 13 ]}\nobreak
\hyperdef{L}{X7EB7CB0B84963812}{}
{\noindent\textcolor{FuncColor}{$\triangleright$\enspace\texttt{CocTestOnClusterList({\mdseries\slshape CL[, opt]})\index{CocTestOnClusterList@\texttt{CocTestOnClusterList}}
\label{CocTestOnClusterList}
}\hfill{\scriptsize (function)}}\\
\textbf{\indent Returns:\ }
A list of pairs each containing the path of the respective cluster and the
subbranch resulting from \texttt{CocSplitRoutine}. 



 Tries to split each cluster in \mbox{\texttt{\mdseries\slshape CL}} via \texttt{CocSplitRoutine} (\ref{CocSplitRoutine}). \mbox{\texttt{\mdseries\slshape CL}} must be a cluster list produced by \texttt{ClusterList} (\ref{ClusterList}). This function accepts a record \mbox{\texttt{\mdseries\slshape opt}} of options with the following content: 
\begin{itemize}
\item  \mbox{\texttt{\mdseries\slshape pack}}: A boolean that determines if \texttt{IdPGroup{\textunderscore}Pack} (\ref{IdPGroupuScorePack}) should be applied to the result. 
\item  \mbox{\texttt{\mdseries\slshape order}}: An integer giving the order of the groups with ids in \mbox{\texttt{\mdseries\slshape CL}}. 
\item  \mbox{\texttt{\mdseries\slshape iterate}}: An integer that determines the maximal number of iterations done while
searching parameters. 
\item  \mbox{\texttt{\mdseries\slshape verbose}}: An boolean that determines if progress should be printed. 
\end{itemize}
 The default value of \mbox{\texttt{\mdseries\slshape opt}} is \mbox{\texttt{\mdseries\slshape rec(pack}}:=\texttt{true}, \mbox{\texttt{\mdseries\slshape order}}:=512, \mbox{\texttt{\mdseries\slshape iterate}}:=1,\mbox{\texttt{\mdseries\slshape verbose}}:=\texttt{true}\mbox{\texttt{\mdseries\slshape )}}. 
\begin{Verbatim}[commandchars=!@|,fontsize=\small,frame=single,label=Example]
  !gapprompt@gap>| !gapinput@tree := TestOnClusterList([[3,5,6,7,8]],"CocList",rec(verbose:=false));;|
  !gapprompt@gap>| !gapinput@CL := ClusterList(tree);|
  [ rec( descl := [  ], ids := [ 3, 8 ], path := [ 337 ], size := 2,
         tests := [ "CocList" ] ),
    rec( descl := [  ], ids := [ 5, 6, 7 ], path := [ 32173 ], size := 3,
         tests := [ "CocList" ] ) ]
  !gapprompt@gap>| !gapinput@TR := CocTestOnClusterList(CL,rec(verbose:=false));;|
  !gapprompt@gap>| !gapinput@AddClusterTestResults(tree, TR);;|
  !gapprompt@gap>| !gapinput@ClusterList(tree);|
  [ rec( descl := [ [ 302004 ] ], ids := [ 3 ], path := [ 337, 1363 ], size := 1,
         tests := [ "CocList" ]),
    rec( descl := [ [ 302004 ] ], ids := [ 8 ], path := [ 337, 46760 ], size := 1,
         tests := [ "CocList" ]),
    rec( descl := [ [ 104003 ] ], ids := [ 5 ], path := [ 32173, 1214 ],
         size := 1, tests := [ "CocList" ] ),
    rec( descl := [ [ 104003 ], [ 40204005 ] ], ids := [ 6 ],
         path := [ 32173, 16, 20792 ], size := 1, tests := [ "CocList" ] ),
    rec( descl := [ [ 104003 ], [ 40204005 ] ], ids := [ 7 ],
         path := [ 32173, 16, 46881 ], size := 1, tests := [ "CocList" ] ) ]
\end{Verbatim}
 }

 

\subsection{\textcolor{Chapter }{ParametersRandPresentTest}}
\logpage{[ 2, 5, 14 ]}\nobreak
\hyperdef{L}{X7FBD2B9C83B1496A}{}
{\noindent\textcolor{FuncColor}{$\triangleright$\enspace\texttt{ParametersRandPresentTest({\mdseries\slshape CL, order, maxl})\index{ParametersRandPresentTest@\texttt{ParametersRandPresentTest}}
\label{ParametersRandPresentTest}
}\hfill{\scriptsize (function)}}\\
\textbf{\indent Returns:\ }
A list consisting of pairs of a list of ids and a record containing the
parameters for a random presentation test via \texttt{ID{\textunderscore}GROUP{\textunderscore}FUNCS[18]}. 



 Determines the needed parameters for the groups of size \mbox{\texttt{\mdseries\slshape order}} in the clusters \mbox{\texttt{\mdseries\slshape cl}} of \mbox{\texttt{\mdseries\slshape CL}}. Creates the clusters of conjugcy classes and alters them accordingly to the \mbox{\texttt{\mdseries\slshape cl.descl}}. Afterwards minimal generators for all groups and their position among the
clusters of conjugacy classes. If the positions do not coincide, the positions
of the generators of the first group are selected and the generators of the
other groups are selected from the same clusters. Otherwise the positions and
generators are sorted by the positions. Sequently the function tries to refine
the corresponding clusters to limit the choices to choose from and repeat the
process until either the number of refinments is equal or larger than \mbox{\texttt{\mdseries\slshape maxl}} or no more refinments could be found. 
\begin{Verbatim}[commandchars=!@|,fontsize=\small,frame=single,label=Example]
  !gapprompt@gap>| !gapinput@ParametersRandPresentTest([rec(descl:=[], ids:=[1367, 1371])], 512, 0);|
  [ [ [ 1367, 1371 ], rec( descl := [  ],
       fp := [ 39080, 80814 ], func := 18, pos := [ 8, 10 ] ) ] ]
  !gapprompt@gap>| !gapinput@ParametersRandPresentTest([rec(descl:=[], ids:=[1367, 1371])], 512, 1);|
  [ [ [ 1367, 1371 ], rec( descl := [ 106003 ],
       fp := [ 39080, 80814 ], func := 18, pos := [ 9, 11 ] ) ] ]
\end{Verbatim}
 }

 }

 }

 \def\indexname{Index\logpage{[ "Ind", 0, 0 ]}
\hyperdef{L}{X83A0356F839C696F}{}
}

\cleardoublepage
\phantomsection
\addcontentsline{toc}{chapter}{Index}


\printindex

\immediate\write\pagenrlog{["Ind", 0, 0], \arabic{page},}
\immediate\write\pagenrlog{["Ind", 0, 0], \arabic{page},}
\immediate\write\pagenrlog{["Ind", 0, 0], \arabic{page},}
\immediate\write\pagenrlog{["Ind", 0, 0], \arabic{page},}
\immediate\write\pagenrlog{["Ind", 0, 0], \arabic{page},}
\immediate\write\pagenrlog{["Ind", 0, 0], \arabic{page},}
\immediate\write\pagenrlog{["Ind", 0, 0], \arabic{page},}
\immediate\write\pagenrlog{["Ind", 0, 0], \arabic{page},}
\immediate\write\pagenrlog{["Ind", 0, 0], \arabic{page},}
\immediate\write\pagenrlog{["Ind", 0, 0], \arabic{page},}
\immediate\write\pagenrlog{["Ind", 0, 0], \arabic{page},}
\immediate\write\pagenrlog{["Ind", 0, 0], \arabic{page},}
\immediate\write\pagenrlog{["Ind", 0, 0], \arabic{page},}
\immediate\write\pagenrlog{["Ind", 0, 0], \arabic{page},}
\immediate\write\pagenrlog{["Ind", 0, 0], \arabic{page},}
\immediate\write\pagenrlog{["Ind", 0, 0], \arabic{page},}
\immediate\write\pagenrlog{["Ind", 0, 0], \arabic{page},}
\immediate\write\pagenrlog{["Ind", 0, 0], \arabic{page},}
\immediate\write\pagenrlog{["Ind", 0, 0], \arabic{page},}
\immediate\write\pagenrlog{["Ind", 0, 0], \arabic{page},}
\immediate\write\pagenrlog{["Ind", 0, 0], \arabic{page},}
\immediate\write\pagenrlog{["Ind", 0, 0], \arabic{page},}
\immediate\write\pagenrlog{["Ind", 0, 0], \arabic{page},}
\immediate\write\pagenrlog{["Ind", 0, 0], \arabic{page},}
\immediate\write\pagenrlog{["Ind", 0, 0], \arabic{page},}
\immediate\write\pagenrlog{["Ind", 0, 0], \arabic{page},}
\immediate\write\pagenrlog{["Ind", 0, 0], \arabic{page},}
\immediate\write\pagenrlog{["Ind", 0, 0], \arabic{page},}
\immediate\write\pagenrlog{["Ind", 0, 0], \arabic{page},}
\immediate\write\pagenrlog{["Ind", 0, 0], \arabic{page},}
\immediate\write\pagenrlog{["Ind", 0, 0], \arabic{page},}
\immediate\write\pagenrlog{["Ind", 0, 0], \arabic{page},}
\immediate\write\pagenrlog{["Ind", 0, 0], \arabic{page},}
\immediate\write\pagenrlog{["Ind", 0, 0], \arabic{page},}
\immediate\write\pagenrlog{["Ind", 0, 0], \arabic{page},}
\immediate\write\pagenrlog{["Ind", 0, 0], \arabic{page},}
\immediate\write\pagenrlog{["Ind", 0, 0], \arabic{page},}
\immediate\write\pagenrlog{["Ind", 0, 0], \arabic{page},}
\immediate\write\pagenrlog{["Ind", 0, 0], \arabic{page},}
\immediate\write\pagenrlog{["Ind", 0, 0], \arabic{page},}
\immediate\write\pagenrlog{["Ind", 0, 0], \arabic{page},}
\immediate\write\pagenrlog{["Ind", 0, 0], \arabic{page},}
\immediate\write\pagenrlog{["Ind", 0, 0], \arabic{page},}
\immediate\write\pagenrlog{["Ind", 0, 0], \arabic{page},}
\immediate\write\pagenrlog{["Ind", 0, 0], \arabic{page},}
\immediate\write\pagenrlog{["Ind", 0, 0], \arabic{page},}
\immediate\write\pagenrlog{["Ind", 0, 0], \arabic{page},}
\immediate\write\pagenrlog{["Ind", 0, 0], \arabic{page},}
\immediate\write\pagenrlog{["Ind", 0, 0], \arabic{page},}
\immediate\write\pagenrlog{["Ind", 0, 0], \arabic{page},}
\immediate\write\pagenrlog{["Ind", 0, 0], \arabic{page},}
\immediate\write\pagenrlog{["Ind", 0, 0], \arabic{page},}
\immediate\write\pagenrlog{["Ind", 0, 0], \arabic{page},}
\immediate\write\pagenrlog{["Ind", 0, 0], \arabic{page},}
\immediate\write\pagenrlog{["Ind", 0, 0], \arabic{page},}
\immediate\write\pagenrlog{["Ind", 0, 0], \arabic{page},}
\immediate\write\pagenrlog{["Ind", 0, 0], \arabic{page},}
\immediate\write\pagenrlog{["Ind", 0, 0], \arabic{page},}
\immediate\write\pagenrlog{["Ind", 0, 0], \arabic{page},}
\immediate\write\pagenrlog{["Ind", 0, 0], \arabic{page},}
\immediate\write\pagenrlog{["Ind", 0, 0], \arabic{page},}
\immediate\write\pagenrlog{["Ind", 0, 0], \arabic{page},}
\immediate\write\pagenrlog{["Ind", 0, 0], \arabic{page},}
\immediate\write\pagenrlog{["Ind", 0, 0], \arabic{page},}
\immediate\write\pagenrlog{["Ind", 0, 0], \arabic{page},}
\immediate\write\pagenrlog{["Ind", 0, 0], \arabic{page},}
\immediate\write\pagenrlog{["Ind", 0, 0], \arabic{page},}
\immediate\write\pagenrlog{["Ind", 0, 0], \arabic{page},}
\immediate\write\pagenrlog{["Ind", 0, 0], \arabic{page},}
\immediate\write\pagenrlog{["Ind", 0, 0], \arabic{page},}
\newpage
\immediate\write\pagenrlog{["End"], \arabic{page}];}
\immediate\closeout\pagenrlog
\end{document}
